\subsection{Modified realizability and total extracted functions}
\label{sec:p07-modified-realizability}

Number realizers may be partial programs that terminate only on
realizing inputs. System-T terms instead denote total functionals on
their finite types. This difference supplies a nonderivability argument
for Markov's principle. The interpretation used here is modified
realizability; it is distinct from both ordinary number realizability
and the proof-carrying relation used for HA proof extraction.

\begin{hbdefinition}{def:p07-modified-realizability}{Modified realizability}
Interpret ground type by \(\mathbb N\) and an arrow type by all
set-theoretic functions between the interpreted types. A tuple below
is a finite metalinguistic list; it adds no product type to System T.
Associate a positive witness tuple type \(\rho(A)\) with each
arithmetic formula: an atom has an empty tuple, conjunction concatenates
the two tuple types, disjunction adds a numerical tag and both tuples,
existential quantification adds a numerical witness, and universal
quantification replaces each component type \(\sigma\) by
\(0\to_{\mathrm{ty}}\sigma\). For implication \(A\to B\),
replace each component \(\tau\) of \(\rho(B)\) by the arrow type
from the successive components of \(\rho(A)\) to \(\tau\).
If the input tuple is empty there are no input arrows; if the output
tuple is empty there are no components at all.

Define the semantic relation \(\operatorname{mr}_A\) recursively:
\begin{enumerate}
\item \(\operatorname{mr}_P()\) means that the numerical atom \(P\)
      is true;
\item \(\operatorname{mr}_{A\land B}(\vec a,\vec b)\) means
      \(\operatorname{mr}_A(\vec a)\land\operatorname{mr}_B(\vec b)\);
\item \(\operatorname{mr}_{A\lor B}(d,\vec a,\vec b)\) means
      \((d=0\land\operatorname{mr}_A(\vec a))\lor
      (d=1\land\operatorname{mr}_B(\vec b))\);
\item \(\operatorname{mr}_{A\to B}(\vec F)\) means that for every
      \(\vec a\) of the input types,
      \(\operatorname{mr}_A(\vec a)\) implies
      \(\operatorname{mr}_B(\vec F(\vec a))\);
\item \(\operatorname{mr}_{\forall nA(n)}(\vec F)\) means
      \(\operatorname{mr}_{A(n)}(\vec F(n))\) for every natural
      number \(n\);
\item \(\operatorname{mr}_{\exists nA(n)}(n,\vec a)\) means
      \(\operatorname{mr}_{A(n)}(\vec a)\).
\end{enumerate}
Object-language numerical parameters remain parameters of this relation.
All quantification over witness tuples in these clauses ranges over
the full interpreted finite types, not just denotations of terms.
In particular negation has an empty positive tuple, and its condition
is that no positive witness realizes its argument.
\end{hbdefinition}

\begin{hblemma}{thm:p07-system-t-numerical-computability}{proved=annex:p7.mr.1}{Numerical computation}
Every closed System-T term of ground type has a finite compatible
beta/recursor reduction to a numeral. Such a numeral can be found
effectively and has the term's unique value in the full finite-type
interpretation. Consequently every closed System-T term of type
\(0\to_{\mathrm{ty}}0\) defines a total computable numerical
function. This assertion does not require a claim that every possible
reduction sequence terminates.
\end{hblemma}

\begin{hbtheorem}{thm:p07-modified-realizability-soundness}{proved=annex:p7.mr.2}{Modified soundness}
An HA proof of \(A(\vec n)\) yields witness terms effectively.
They have the component types \(\rho(A)\), with free parameters
among \(\vec n\). For every numerical tuple their denotations satisfy
\(\operatorname{mr}_{A(\vec n)}\). Closed conclusions yield closed
terms. No constants beyond System T are used.
\end{hbtheorem}

\begin{hblemma}{thm:p07-bounded-modified-realizability}{proved=annex:p7.mr.3}{Bounded truth}
For every fixed bounded arithmetic formula \(B(\vec n)\):
\begin{enumerate}
\item any tuple satisfying \(\operatorname{mr}_B\) implies the
      numerical truth of \(B\);
\item System T defines a canonical tuple \(\vec b(\vec n)\),
      total for all parameters, which satisfies
      \(\operatorname{mr}_{B(\vec n)}\) whenever \(B(\vec n)\)
      is true;
\item System T defines a numerical decision term for its truth.
\end{enumerate}
Thus truth is equivalent to existence of a modified realizer for these
bounded formulas. In particular there is a uniform System-T witness
for \(\forall t(B(\vec n,t)\lor\neg B(\vec n,t))\).
\end{hblemma}
